%---------------------------------------------------------------------------%
%-> Chinese abstract content for CombinedThesis
%---------------------------------------------------------------------------%
同步定位与地图构建（Simultaneous Localization and Mapping, SLAM）是自主移动机器人、无人机及增强现实设备实现环境感知与自主导航的核心技术。视觉 SLAM 后端优化涉及大规模稀疏矩阵构建与线性方程组求解，具有极高的计算复杂度与内存访问压力，传统嵌入式 ARM 处理器难以实时完成。针对这一核心矛盾，本文提出并系统实现了一套基于 FPGA 的软硬协同 SLAM 后端加速方案，将"舒尔消元 + EKF 增量更新"的算法策略映射为可在 Zynq UltraScale+ 平台高效运行的异构硬件架构。

在算法层面，本文以 SchurVINS 框架为基础，利用光束法平差（BA）中 Hessian 矩阵的箭形块稀疏结构，通过舒尔补消元将路标点状态从系统中消去，再以单次 EKF 观测更新替代传统 Levenberg--Marquardt 多轮迭代求解，将后端计算主体归结为雅可比累加与 $3 \times 3$ 矩阵求逆两个规则化计算核，实现了算法复杂度与硬件友好性的统一。

在硬件设计层面，本文完成了三项核心工作：（1）设计了 5 级深度流水线雅可比计算模块，覆盖从世界坐标变换、重投影残差计算到 Hessian 矩阵在线累加的完整链路，并通过切换外参常量实现左右目双相机的单硬件路径复用；（2）提出了基于 8 位稀疏观测位图（SOB）编码的两级动态门控机制，在流水线前端实现状态级旁路与相机级时钟门控，自动跳过零块计算，在典型稀疏率 50\% 场景下将总浮点运算量降低约 45\%；（3）实现了针对 $3 \times 3$ 对称矩阵的并行化 Schur 消元核心，以代数余子式法并行计算逆矩阵并利用对称性将输出写操作减少约 50\%，Hessian 矩阵片上 BRAM 占用仅为 67~KB。

在系统架构层面，本文遵循逻辑密集型任务归 PS、数据密集型规则计算归 PL 的原则，在 PS 端 ARM 处理器上运行 SVO 2.0 视觉前端追踪、滑动窗口状态管理与 EKF 位姿更新，将雅可比生成、Hessian 累加与 Schur 消元完全卸载至 PL 端 FPGA 流水线。三通道 AXI 通信架构与双缓冲乒乓调度策略有效隐藏了片外存储访问延迟。

实验结果表明，在 Zynq UltraScale+ ZCU102 平台上，Schur Kernel 将单次后端优化延迟从 2.84~ms 压缩至 0.76~ms，实现 3.74 倍加速比，系统后端优化频率稳定在 25~Hz 以上。在轨迹精度方面，EuRoC 数据集上 FPGA 方案平均绝对轨迹误差（ATE）为 0.073~m，与全软件基准的 0.075~m 精度等价；TUM-VI 数据集上 FPGA 方案平均 ATE 为 0.152~m，优于软件基准的 0.182~m。在能效方面，单帧能耗仅为 19.1~mJ，BRAM 占用为 67~KB，均为同类方案中最优。户外真实场景的部署测试进一步验证了系统的工程可行性。

\keywords{视觉 SLAM，视觉惯性里程计，FPGA，硬件加速，软硬协同，舒尔补消元，稀疏观测位图}
